\documentclass{Classes/PhDThesisPSnPDF}
\usepackage[utf8]{inputenc}
\usepackage[T1]{fontenc}
\usepackage{geometry}
\usepackage{graphicx}
\usepackage{float}
\usepackage{hyperref}
\usepackage[english,indonesian]{babel}

\title{Optimasi Pemrosesan Template LaTeX Menggunakan Pendekatan Berbasis AI}
\author{Academic Support Agent}
\date{\today}

\begin{document}
\maketitle

\begin{abstract}

\end{abstract}

\chapter{PENDAHULUAN}


\chapter{TINJAUAN PUSTAKA}


% Original content (commented out):
% % ======================================
%  Ph.D. THESIS TEMPLATE
% --------------------------------------
%  Ph.D. in Economics
% --------------------------------------
%  Università degli Studi Roma Tre
%  Facoltà di Economia "Federico Caffè"
% ======================================

% The following LaTeX template has been developed based on the one provided by the University of Cambridge, available at the following link: http://www-h.eng.cam.ac.uk/help/tpl/textprocessing/ThesisStyle/

% Inquires and suggestions can be submitted to Nicola Caravaggio
% nicola.caravaggio.cilu@gmail.com

% Note:
% Several commands have to be commented out to be activated!

% ============================================================================

\documentclass[article,12pt,times,print,authoryear,titlepage,index]{Classes/PhDThesisPSnPDF}

% ==========
%  Packages
% ==========
\usepackage[T1]{fontenc}
\usepackage[utf8]{inputenc} 
\usepackage{lipsum} 
\usepackage{natbib}
\usepackage{frontespizio}
\usepackage{lipsum} % Just to create text examples
\usepackage[inlinechap]{settings/fncystyle} % Chapter style (Baskerville as deafult font)

% Fonts
%\usepackage{mathpazo} % Baskerville
\usepackage{times} % Times New Roman
%\usepackage{ebgaramond} % EB Garamond
%\usepackage{garamondx} % Garamond X
%\usepackage{CormorantGaramond} % Cormorant Garamond

% Dictionary
\usepackage[italian,english]{babel}

% Tools for charts
\usepackage{array}
\usepackage{siunitx} 
\usepackage{graphicx}
\usepackage{multirow} 
\usepackage[normalem]{ulem} 
\usepackage{booktabs,caption} 
\usepackage[flushleft]{threeparttable} % For tables' notes
\usepackage{longtable} % Splits long tables on several pages
\usepackage{ltxtable}
\usepackage{adjustbox}
\usepackage{rotating}
\usepackage{subcaption}
\usepackage{framed}
\usepackage{varwidth}
\usepackage{lscape}
\renewcommand{\floatpagefraction}{.9} % To set figures' floating
\usepackage[font=small,skip=2pt]{caption} % To reduce title space
\newcommand{\rr}{\raggedright}
\newcommand{\tn}{\tabularnewline}
\usepackage[capitalise]{cleveref} % Capitalize multiple references
\usepackage{hhline} % Double line in charts
\newcommand\T{\rule{0pt}{2.6ex}} % Top strut (type \T for top space)
\newcommand\B{\rule[-1.2ex]{0pt}{0pt}} % Bottom strut (type \B for bottom space)

% Tools for equations
\usepackage{mathtools}
\usepackage{amsfonts}
\usepackage{array,amsmath} 
\usepackage{tabularx}

% Enumeration
\usepackage{enumitem} 
\setlist{parsep=0pt,listparindent=\parindent}

% Inspirational quotes
\usepackage{epigraph} % For inspirational quotes
\setlength\epigraphwidth{6cm} % For inspirational quotes
%\setlength\epigraphrule{0pt} % Inspirational quote without line

% References
\usepackage{csquotes}
\usepackage{url} % Add link to references
\usepackage{hyperref} % Add link to references
\usepackage{hyperref} % To make references among thesis' parts

% Customize pages' strucutre
\usepackage{geometry}
\usepackage{fancyhdr} 

% Other packages
\usepackage{ulem} % Different kind of underlined text
\usepackage{afterpage} % To break pages and start in the next one
\usepackage{sepfootnotes} % Useful for long footnotes (\sepfootnotecontent)
\usepackage{mhchem} % Useful for subscript in text (es. \ce{CO2})
\usepackage{csquotes} % For quotes
\usepackage{epstopdf} % To convert .eps (STATA) files into .pdf
\usepackage{setspace} % To set line spacing
\usepackage{ragged2e} % Ragged text
\usepackage{color} % To change text color
\usepackage{ragged2e} % Justification
\usepackage{gensymb} % Symbols
\usepackage[space]{grffile} % Change directory
\usepackage[en-US]{datetime2} % Date
\usepackage{relsize} % In text subscript and superscript
%\usepackage{datetime} % Date
%\usepackage{lettrine} % Start paragraph with a big letter
%\renewcommand{\LettrineTextFont}{\normalfont} % No bold font after lettrine
%\usepackage{titling} % Pro­vides con­trol over the type­set­ting of the \maketi­tle
%\usepackage{pdfpages} % Extract only some pages fromt he document

\def\textsubscript#1{\ensuremath{_{\mbox{\textscale{.6}{#1}}}}}

\useunder{\uline}{\ul}{} % ? (do not unselect it!)

% Nomenclature
\usepackage{nomencl} % List of Abbreviations
\makenomenclature
% To divide the nomenclature into different sub-categories (e.g. A,B,C) you have to comment-out the following commands.

%\usepackage{etoolbox}
%\renewcommand\nomgroup[1]{%
%	\item[\bfseries
%	\ifstrequal{#1}{A}{}{%
%	\ifstrequal{#1}{B}{}{%
%	\ifstrequal{#1}{C}{}{}}}%
%]}

% To generate various nomenclatures use following command in your document where "prefix" is the sub-category selectes (e.g. A,B,C).

%\nomenclature[prefix]{symbol}{description}

% To change the name of the nomenclature section check the "Preamble" file.

% In order to have the Nomenclature shwed in the document you have to copy and paste the following text
% ---> makeindex %.nlo -s nomencl.ist -o %.nls -t %.nlg
% In the commands of your editor software in the section "MakeIndex"
% For example, with TeXstudio you have to go: 
% Options/Configure TeXstudio.../Commands, then in the space of "MakeIndex"
% For example, with Texmaker you have to go: 
% Options/Configure Texmaker/Commands, then in the space of "Makeindex"
% Therefore, you have to compile your document in three steps:
% 1) Compile the .pdf
% 2) Compile MakeIndex (with TeXstudio you have to go to Tools/Index)
% 3) Compile again the .pdf, then show it!

% N.B.
% Additional packages are listed in the "Preamble" file.

% ============================================================================

% ==========
%  Preamble
% ==========
% Contains packages and user-defined commands and settings
\input{Preamble/preamble}


% ===================
%  Abstract Separate 
% ===================
% To printout only the titlepage and the abstract with the PhD title and the author name 
% use the `abstract' option in the document class

\ifdefineAbstract
\pagestyle{empty}
\includeonly{Declaration/declaration, Abstract/abstract}
\fi


% ==============
%  Chapter Mode 
% ==============
% The chapter mode allows user to only print particular chapters with references
% Title, Contents, Frontmatter are disabled by default
% Useful option to review a particular chapter or to send it to your supervisor(s)
% To use choose `chapter' option in the document class

%\ifdefineChapter
% \includeonly{Chapter1/chapter1}
% \includeonly{Chapter2/chapter2}
% \includeonly{Chapter3/chapter3}
% \includeonly{Chapter4/chapter4}
% \includeonly{Chapter5/chapter5}
% \includeonly{Chapter6/chapter6}
%\fi


% ==============
%  Front Matter 
% ==============
\begin{document}
	
\pagestyle{empty} % To delete numbers at the bottom of the page
\include{firstpage} % First page 	
\include{frontispiece} % Frontispiece
	
%\maketitle
\frontmatter
%\include{Dedication/dedication}
%\include{Quote/quote}
%\include{Acknowledgement/acknowledgement}
\include{Abstract/abstract}
	
	
% ================================
%  Adding TOC and List of Figures 
% ================================

\pagestyle{fancy}	
\tableofcontents
\listoffigures	
\listoftables
% To print out the nomenclature comment-out the following command
\printnomenclature[4em] % 4em is the space between symbols and descriptions

	
% ========================
%  Main Matter - Chapters
% ========================
\mainmatter
	
\include{Introduction/introduction}
\include{Chapter1/chapter1}
\include{Chapter2/chapter2}
\include{Chapter3/chapter3}
%\include{Chapter4/chapter4}
%\include{Chapter5/chapter5}
%\include{Chapter6/chapter6}
\include{Conclusions/conclusions}	
	
	
% =============
%  Back Matter 
% =============
% Backmatter should be commented out, if you are using appendices after References and if you do not want "A", "B", etc. for appendices.
%\backmatter
	
	
% ============
%  Appendices
% ============
\begin{appendices} % Using appendices environment for more functunality
	\renewcommand\appendixname{Appendix}
	%\renewcommand\appendixname{\empty} % Without "Appendix", only letters 
	\include{Appendix1/appendix1}
	%\include{Appendix2/appendix2}
	%\include{Appendix3/appendix3}
	
\end{appendices}
	
	
% ==============
%  Bibliography 
% ==============
\begin{spacing}{0.9}
		
	% To use the conventional natbib style referencing
	% Bibliography style previews: http://nodonn.tipido.net/bibstyle.php
	% Reference styles: http://sites.stat.psu.edu/~surajit/present/bib.htm
		
	\bibliographystyle{apalike}
	%\bibliographystyle{unsrt} % Use for unsorted references  
	%\bibliographystyle{plainnat} % Use this to have URLs listed in References
	\cleardoublepage
	\bibliography{References/references} % Path to your References.bib file
	\renewcommand{\bibname}{References}
		
\end{spacing}

% =======
%  Index 
% =======
\printthesisindex % If index is present
	
% To extract only some specific pages
%\includepdf[pages={1,9-14,18-62}]{thesis.pdf}

\end{document}

\end{document}